\section{Discussion}

The procedure a user follows has four steps. Choose the resolution $r$ that matches the question: a coarse $r$ for ``are the main kinds of post represented'', a fine $r$ for ``has every specific claim been seen''. Embed a pilot of about 2{,}000 documents and estimate $d$ by TwoNN. Fit Equation~\ref{eq:collapse} on nested subsamples of the pilot and check its $R^2$. Read $n^{*}$ from Equation~\ref{eq:nstar} and, after the sample is drawn, report $G\,(1 - n/N)$ as the share of the corpus the sample leaves uncovered. No step needs labels, and the whole runs on a laptop once the embeddings exist.

The survey analogy is exact in two places and instructive in a third. The finite population correction is the same factor for the same reason. The dimension $d$ plays the role of the variance: it is the one property of the population that the sample size formula needs, and like the variance it is estimated from a pilot. The third place is the margin of error. A survey's margin bounds the error of a proportion, and a corpus sample's plain mean already does that job well, to within 0.01 for cluster proportions at $n = 1{,}000$ in our data. The missing mass bounds something else, the mass of content the sample has no example of, which is the error that matters for enumeration: listing the themes, finding the rare claim, checking that a summary had the chance to mention something. The two margins answer different questions, and a report of a corpus sample should carry both.

The theme result explains two failures that an applied reader will have met. A free summary written from a sample is stable across draws at a rate that depends on the summarizer's own granularity, and a theme inventory of fixed size covers at the same rate in every corpus. Both readouts are relative to the corpus, and both are blind to how broad the corpus is at any absolute scale. The breadth that the sample size formula is about lives at a resolution, and only a readout at that resolution can see it.

The collapse diagnostic has a use beyond the formula. A corpus whose isolation surface does not collapse is one where slogans and discussion, or several registers, sit at different scales, and a single sample size will oversample one and undersample the other. The probe of local dimension across clusters locates the registers, and a stratified sample with a size per stratum is the natural next step.
